\section{Orthography}
\subsection{Consonants}
\begin{table}[H]
  \caption{\label{tab:orthography-consonants}Consonant orthography}
  \renewcommand\arraystretch{2}
  \setlength\extrarowheight{-2pt}
  \begin{tabular}{|c|c|l|}
    \hline
    \rowcolor{greyhead}
    Allophone & Grapheme & Description \\
    \hline
    \dsil [p] & \dsil ‹p› & Voiceless bilabial plosive \\
    \hline
    \rowcolor{greyrow}
    \dsil [t̪] & \dsil ‹t› &  Voiceless laminal dental plosive \\
    \hline
    \dsil [m] & \dsil ‹m› & Voiced bilabial nasal \\
    \hline
    \rowcolor{greyrow}
    \dsil [n̪] & \dsil ‹n› & Voiced laminal dental nasal \\
    \hline
    \dsil [r] & \dsil ‹r› & Voiced alveolar trill \\
    \hline
    \rowcolor{greyrow}
    \dsil [ɾ] & \dsil ‹r› & Voiced alveolar tap or flap \\
    \hline
    \dsil [ɸ] & \dsil ‹ph› & Voiceless bilabial fricative \\
    \hline
    \rowcolor{greyrow}
    \dsil [ʍ] & \dsil ‹wh› & Voiceless labialized velar approximant \\
    \hline
    \dsil [θ] & \dsil ‹th› & Voiceless dental fricative \\
    \hline
    \rowcolor{greyrow}
    \dsil [s] & \dsil ‹s› & Voiceless alveolar fricative \\
    \hline
    \dsil [ʃ] & \dsil ‹sh› & Voiceless postalveolar fricative \\
    \hline
    \rowcolor{greyrow}
    \dsil [h] & \dsil ‹h› & Voiceless glottal fricative \\
    \hline
    \dsil [w] & \dsil ‹w› & Voiced labialized velar approximant \\
    \hline
    \rowcolor{greyrow}
    \dsil [l] & \dsil ‹l› & Voiced alveolar lateral approximant \\
    \hline
  \end{tabular}
\end{table}
\subsection{Vowels}
\begin{table}[H]
  \caption{\label{tab:orthography-vowels}Vowel orthography}
  \renewcommand\arraystretch{2}
  \setlength\extrarowheight{-2pt}
  \begin{tabular}{|c|c|l|}
    \hline
    \rowcolor{greyhead}
    Allophone & Grapheme & Description \\
    \hline
    \dsil [i] & \dsil ‹i› &  Close front unrounded \\
    \hline
    \rowcolor{greyrow}
    \dsil [u] & \dsil ‹u› &  Close back rounded \\
    \hline
    \dsil [e̞] & \dsil ‹e› & Mid front unrounded \\
    \hline
    \rowcolor{greyrow}
    \dsil [o̞] & \dsil ‹o› & Mid back rounded \\
    \hline
    \dsil [ä] & \dsil ‹a› & Open central unrounded \\
    \hline
  \end{tabular}
\end{table}
\clearpage
\subsection{Vowel Pairs}

Vowel pairs are never diphthongs in Phi; they are always pronounced separately.
The vowel pairs specified in Table~\ref{tab:orthography-pairs} show the vowels
in hiatus.

\begin{table}[H]
  \caption{\label{tab:orthography-pairs}Vowel pair orthography}
  \renewcommand\arraystretch{2}
  \setlength\extrarowheight{-2pt}
  \begin{tabular}{|c|c|l|}
    \hline
    \rowcolor{greyhead}
    Allophone & Grapheme \\
    \hline
    \dsil [i.e̞] & \dsil ‹ie› \\
    \hline
    \rowcolor{greyrow}
    \dsil [i.o̞] & \dsil ‹io› \\
    \hline
    \dsil [i.ä] & \dsil ‹ia› \\
    \hline
    \rowcolor{greyrow}
    \dsil [e̞.i] & \dsil ‹ei› \\
    \hline
    \dsil [e̞.o] & \dsil ‹eo› \\
    \hline
    \rowcolor{greyrow}
    \dsil [e̞.ä] & \dsil ‹ea› \\
    \hline
    \dsil [o̞.i] & \dsil ‹oi› \\
    \hline
    \rowcolor{greyrow}
    \dsil [o̞.e̞] & \dsil ‹oe› \\
    \hline
    \dsil [o̞.ä] & \dsil ‹oa› \\
    \hline
    \rowcolor{greyrow}
    \dsil [ä.i] & \dsil ‹ai› \\
    \hline
    \dsil [ä.e̞] & \dsil ‹ae› \\
    \hline
    \rowcolor{greyrow}
    \dsil [ä.o̞] & \dsil ‹ao› \\
    \hline
  \end{tabular}
\end{table}
